\section{Introducción}
\label{Introduccion}

En los últimos años, los drones han emergido como herramientas clave en los conflictos bélicos modernos, representando una evolución significativa en la tecnología militar. Particularmente en el contexto de la guerra Ruso-Ucraniana, estos dispositivos han sido adoptados ampliamente como armas de guerra, destacándose por su precisión, versatilidad y relativo bajo costo en comparación con los sistemas tradicionales. Sin embargo, esta innovación tecnológica también ha expuesto vulnerabilidades en los sistemas de defensa actuales, los cuales no han logrado ser completamente efectivos frente a esta nueva amenaza.

La principal debilidad de las defensas convencionales radica en su incapacidad para contrarrestar la precisión y la adaptabilidad de los drones, así como en el elevado costo asociado al despliegue y mantenimiento de sistemas de misiles antiaéreos. En este contexto, surge la necesidad imperiosa de desarrollar soluciones innovadoras que permitan identificar y neutralizar estas amenazas de manera eficiente y económica. 

El presente trabajo aborda esta problemática mediante la integración de técnicas avanzadas de visión computacional y aprendizaje automático en un robot militar, diseñado para operar en entornos hostiles. Este enfoque busca optimizar la detección y eliminación de objetivos, ofreciendo una alternativa robusta y efectiva a los sistemas de defensa tradicionales.

